\section{A direct quartic in 780 unrestricted integers}\label{dc:section}

The second construction reduces the direct $C_{222}$ quartic of
\cite{direct407} from 888 to 780 integer unknowns.  It keeps degree four,
reduces the collected support from 10,156 to 9,628 monomials, and leaves
the maximum absolute expanded coefficient unchanged at
$2\,445\,790\,974\,704$.  Its new ingredient is an order of proof:
the signs of the private power quotients are recovered \emph{after}
the power graphs have been shown correct.  This removes independent
sign guards without weakening the final arithmetic specification.

\subsection{Exact object and inherited source}

Let $r_0,\ldots,r_{433}$ be the literal quadratic rows in this component's
\texttt{equations/equation.json}, using its listed order of 780 variables.
Define the fixed integer polynomial
\begin{equation}\label{dc:polynomial}
 P_{780}(x_0,\ldots,x_{779})=\sum_{j=0}^{433}r_j(x)^2.
\end{equation}
All 780 variables range over $\Z$.  The complete collected polynomial is
also supplied as \texttt{expanded\_polynomial.json}; this is an ordinary
polynomial with literal integer coefficients and fixed exponents.
The coefficient of $x_0^4$ is 3, so its degree is exactly four.

The pinned immediate source is the archive
\texttt{ConZFC\_407\_888\_Quartics\_2026-10-24.zip}
at ProveIt commit
\texttt{ec80841c07e91be8956ce6ed4af9a9c0b782d3d9}.
It is retained byte for byte, with its original notices; the separate
\texttt{SOURCE\_PINS.json} records the archive digest and the exact members
used.  Write $S(y)=0$ for its 606 quadratic source rows on $y\in\N^{690}$.
The inherited arithmetic and logical calibration is
\begin{equation}\label{dc:source}
 \bigl(\exists y\in\N^{690}: S(y)=0\bigr)
 \ \longleftrightarrow\ C_{222}(1)
 \ \longleftrightarrow\
 \exists p\;\operatorname{Prf}_{\mathrm{ZFC}}(p,\bot).
\end{equation}
The source report identifies the inherited calibration theorems
\begin{quote}\small\ttfamily
KEndpoint.c222\_iff\_zf\_refutation\\
ConstructibleChoice.zf\_refutation\_iff\_zfc\_refutation
\end{quote}
Their calibration snapshot is
\begin{center}\small
\texttt{08200eb10fde2db408c69dffd840d6920a569b1f}.
\end{center}
Here we prove the new existence equivalence with $S$; the proof of
\eqref{dc:source} is inherited from that explicitly identified source.

\subsection{Guarded signs and a staged inference certificate}

For a selected source coordinate $y_v$, three private integer unknowns
enforce the quadratic guard
\begin{equation}\label{dc:guard}
 y_v=a_v^2+b_v^2+c_v(c_v+1).
\end{equation}
This forces $y_v\ge0$ for all integers $a_v,b_v,c_v$.  Conversely, every
nonnegative integer has such a representation: apply Legendre's
three-square theorem to $4y_v+1$, whose residue modulo four forces one
of the three squares to be odd and the other two even, then subtract
one and divide by four.  The parent used 167 such guards.  We use 131.

Ordinary sign propagation is performed on the unchanged source rows.
An oriented row is factored exactly as
$a(\widehat y)y_v-b(\widehat y)=0$, with neither $a$ nor $b$ involving
$y_v$.  Previously proved lower bounds $a\ge1$ and $b\ge0$ imply
$y_v\ge0$; if $b\ge1$, integrality implies $y_v\ge1$.
The checker establishes these bounds from sums of nonnegative
monomials and exact affine squares.  Odd powers need an established
nonnegative factor; even powers need no sign premise.  A monomial
whose factors are all at least one contributes its coefficient to
the constant lower bound.  This last rule can certify a polynomial
with a negative constant.  Every premise must have been established
earlier in the certificate.

Starting from the 131 guards, these elementary inferences prove 623
of the 690 source coordinates nonnegative.  They also prove that every
power base and every exponent is at least one, and prove all sign
premises for the frame below.  In particular, a factorial output used
as a power base receives its strict sign from literal arithmetic rows;
the proof does not first appeal to factorial correctness.

\subsection{The frame bounds do not require power-quotient signs}

There are 23 power calls $(n,k,m)$, nine distinct bases, and 47 distinct
raw atoms occurring as a base, exponent, or proposed output.  Denote
that set of atoms by $\mathcal A$.  The exact shared frame is
\begin{equation}\label{dc:frame}
 W=2+\sum_{a\in\mathcal A}a,
 \qquad Z=z+1,
 \qquad A^2-W^3(W+2)Z^2=1.
\end{equation}
At this point all atoms in $\mathcal A$ are known nonnegative,
$n,k\ge1$, $A\ge0$ and $Z\ge1$.  Consequently
$W\ge2$, $1\le n,k\le W$, $0\le m\le W$, and $A>W^2$.

For $a\ge2$, define the Pell coordinates by
\begin{align}
 X_0(a)&=1,&Y_0(a)&=0,\nonumber\\
 X_{j+1}(a)&=aX_j(a)+(a^2-1)Y_j(a),&
 Y_{j+1}(a)&=X_j(a)+aY_j(a).\label{dc:pell}
\end{align}
The norm in \eqref{dc:frame} is also
$A^2-((W+1)^2-1)(WZ)^2=1$.
Pell-solution classification therefore gives
$A=X_j(W+1)$ and $WZ=Y_j(W+1)$ for some $j>0$.
Reducing \eqref{dc:pell} modulo $W$ gives
$X_j(W+1)\equiv1$ and $Y_j(W+1)\equiv j\pmod W$.
Thus $W\mid j$ and $j\ge W$, whence
\begin{equation}\label{dc:exponentialbound}
 A\ge(W+1)^j\ge(W+1)^W>W^W\ge n^k.
\end{equation}

For each distinct base the literal system defines
\begin{equation}\label{dc:modulus}
 T_n=2An-n^2-1,
 \qquad D=A^2-1.
\end{equation}
Their signs have not been independently guarded.  Since
$1\le n\le W<A$, the function $2An-n^2-1$ is increasing on this
integer interval.  Therefore
\begin{equation}\label{dc:modulusbound}
 T_n\ge2A-2>W^W\ge n^k,
 \qquad T_n>W\ge m,
 \qquad D\ge1.
\end{equation}
We have recovered the nine modulus signs and the discriminant sign
using only the frame and earlier elementary facts.  Elementary
propagation from these facts raises the number of nonnegative source
coordinates to 644.  No private power quotient was used in this step.

\subsection{Index certification precedes power recovery}

Fifteen seed exponents have individual witnesses
$X_i,Y_i,S_i,R_i,e_i,s_i,r_i$, with a shared $U,V,B$ and auxiliary
quotients $b,v,q$.  The exact equations are
\begin{align}
 Y_i&=k_i+e_i,& X_i^2-(A^2-1)Y_i^2&=1,\nonumber\\
 S_i^2-(B^2-1)R_i^2&=1,& S_i&=X_i+Us_i,\nonumber\\
 R_i&=k_i+4Y_i r_i,& U^2-(A^2-1)V^2&=1,\label{dc:index}\\
 B&=A+Ub,&
 V&=\left(\prod_{i=1}^{15}Y_i^2\right)(v+1),\nonumber\\
 B&=1+4V(q+1).&&\nonumber
\end{align}
Every displayed witness has already been proved nonnegative, and
$k_i\ge1$.  Hence $Y_i\ge1$, $V\ge1$, $B\ge5$,
$Y_i^2\mid V$, and $4Y_i\mid(B-1)$.
For each $i$ these equations satisfy the hypotheses of the inherited
Matiyasevich Pell-index criterion
\texttt{PellRepresentation.criterion\_iff\_pell}.  Thus
\begin{equation}\label{dc:indexidentified}
 X_i=X_{k_i}(A),\qquad Y_i=Y_{k_i}(A).
\end{equation}
The three remaining requested pairs are derived by shifts $s=1$ or
$s=2$, using the exact polynomial identities
\begin{equation}\label{dc:addition}
 X_{k+s}=X_kX_s+DY_kY_s,
 \qquad Y_{k+s}=X_kY_s+Y_kX_s,
\end{equation}
where $(X_1,Y_1)=(A,1)$ and $(X_2,Y_2)=(2D+1,2A)$.
Their gate identities and exponent shifts are checked literally.
All required Pell pairs are therefore identified before using the
sign of any power quotient.

\subsection{The quotient sign follows from the exact power}

\begin{lemma}[Delayed power-quotient sign]\label{dc:quotientlemma}
Suppose $A\ge2$, $n,k\ge1$, $T=2An-n^2-1>0$,
$0\le m<T$ and $n^k<T$.  Suppose further that
$X=X_k(A)$, $Y=Y_k(A)$, and that an integer $h$ satisfies
\begin{equation}\label{dc:power}
 X+nY=AY+m+Th.
\end{equation}
Then $m=n^k$ and
\begin{equation}\label{dc:hformula}
 h=\sum_{j=0}^{k-1}n^{k-1-j}Y_j(A)\ge0.
\end{equation}
\end{lemma}

\begin{proof}
Put $H_j=X_j(A)-(A-n)Y_j(A)$.  Direct substitution into
\eqref{dc:pell} gives
\[
 H_0=1,\qquad H_{j+1}=nH_j+TY_j(A).
\]
Induction yields
\begin{equation}\label{dc:Hformula}
 H_k=n^k+T\sum_{j=0}^{k-1}n^{k-1-j}Y_j(A).
\end{equation}
Equations \eqref{dc:power} and \eqref{dc:Hformula} show
$m\equiv n^k\pmod T$, independently of the sign of $h$.
Both numbers lie in $[0,T)$, so $m=n^k$.
Substitution and cancellation of $T\ne0$ then give
\eqref{dc:hformula}.  All its summands are nonnegative.
\end{proof}

The exact recurrence already occurs in the inherited positive-power
construction.  The new use is to apply its soundness argument with
$h$ unrestricted and then recover its sign.  The preceding frame and
index analysis verifies every hypothesis of Lemma~\ref{dc:quotientlemma}
for all 23 calls.  We obtain 23 quotient signs; a final elementary pass
obtains the remaining 23 product signs.  All 690 source coordinates
are now nonnegative.

\subsection{Literal interface and absence of circularity}

The semantic interface is checked against the immutable raw gate ledger,
which contains 1,301 coordinates and 1,217 rows.  Source elimination
replaced 611 of these rows by zero identities; the other 606 become
exactly the literal rows of $S$.  The independent checker reconstructs
every omitted raw coordinate by its monic gate definition and verifies
all these polynomial identities.  This reconstruction is valid at
integer zeros, before any nonnegativity assumption on the restored
coordinates.

The interface file then binds \eqref{dc:frame}, \eqref{dc:modulus},
every instance of \eqref{dc:index}, the three addition formulas, and all
23 equations \eqref{dc:power} to exact raw rows and atoms.  Some named
atoms are related by arithmetic gates---for example, both $t$ and
$t+2$ occur in the distinct-atom sum.  Identity comparison preserves
those gate equations; it does not incorrectly declare such atoms
algebraically independent.

The actual inference order, with checked counts, is
\[
\begin{array}{c|c}
\text{Stage}&\text{Source coordinates known nonnegative}\\\hline
131\text{ private guards and elementary propagation}&623\\
\text{frame bounds and further elementary propagation}&644\\
\text{index identification, power recovery, final propagation}&690.
\end{array}
\]
All index premises are checked before invoking the power lemma.
Power outputs enter $W$ as already nonnegative proposed values; no
claim that they are correct is used to derive the frame bounds.
This establishes the sign argument in an ordered form.

\subsection{Integer projection and the resulting theorem}

Let $G_{131}$ be the 606 source rows together with the 131 guards
\eqref{dc:guard}, on 1,083 unrestricted integer coordinates.
The sign proof shows that every integer zero of $G_{131}$ projects
to a natural zero of $S$.  Conversely any natural zero of $S$ extends
by 131 independent three-square representations.  Therefore
\begin{equation}\label{dc:guardedequivalence}
 \exists y\in\N^{690}\ S(y)=0
 \quad\longleftrightarrow\quad
 \exists z\in\Z^{1083}\ G_{131}(z)=0.
\end{equation}

Only after this conversion do we perform signed coordinate elimination.
A row $z_v-f=0$, or its negative, with $f$ an integer polynomial not
involving $z_v$, permits exact substitution everywhere.  Each reduced
integer zero restores uniquely by $z_v=f$; positivity and divisibility
conditions are unnecessary.  The independent checker expands every
substitution and checks that the remaining rows are quadratic.
There are 303 eliminations: 188 affine and 115 quadratic.  No
unconstrained coordinates are discarded.  Reverse traversal restores
the entire guarded solution uniquely, leaving exactly 780 coordinates
and 434 residuals.  Thus
\begin{equation}\label{dc:projectedequivalence}
 \exists z\in\Z^{1083}\ G_{131}(z)=0
 \quad\longleftrightarrow\quad
 \exists x\in\Z^{780}\ P_{780}(x)=0.
\end{equation}
Uniqueness concerns the monic restoration; the initial three-square
extensions need not be unique.

\begin{theorem}\label{dc:main}
Using the inherited calibration \eqref{dc:source}, the explicit
polynomial \eqref{dc:polynomial} satisfies
\[
 \mathrm{ZFC}\vdash
 \bigl(\nexists x\in\Z^{780}:P_{780}(x)=0\bigr)
 \longleftrightarrow\Con(\mathrm{ZFC}).
\]
\end{theorem}
\begin{proof}
Compose \eqref{dc:source}, \eqref{dc:guardedequivalence}, and
\eqref{dc:projectedequivalence}, using that an integer sum of squares
is zero exactly when every summand is zero.  Negate the resulting
existence equivalence.  The new reasoning is internalizable in ZFC
using integer arithmetic, natural induction, standard Pell theory,
and the three-square theorem; it assumes no consistency statement.
\end{proof}

The 36 guards removed from the immediate parent are the 23 private
power quotients, nine moduli $T_n$, the discriminant $D$, and three
coordinates handled by the stronger positive-factor rule:
\texttt{endpoint.ut}, \texttt{factorial0.output}, and
\texttt{crt.nonnegative\_quotient}.  Each guard costs three private
integer coordinates, so the count decreases by $3\cdot36=108$.

\begin{center}
\begin{tabular}{lrr}
\hline
Quantity&Parent&This construction\\\hline
Independent guards&167&131\\
Integer variables&888&780\\
Quadratic residuals&470&434\\
Collected monomials&10,156&9,628\\
Exact degree&4&4\\
Maximum absolute coefficient&$2\,445\,790\,974\,704$&$2\,445\,790\,974\,704$\\
\hline
\end{tabular}
\end{center}
The variable decrease is 12.16\% within this direct quartic family.
Other repository constructions have fewer variables at different
degrees or coefficient costs.  This result makes no claim of global
minimality or worldwide priority.

The independent finite checker imports neither the producer nor its
elimination module.  It verifies source bytes, the staged sign facts,
all semantic interfaces, every signed inverse substitution, and the
complete collected polynomial.  Separate controls require deliberate
coefficient, index, guard and restoration corruptions to be rejected;
the build is reproduced byte for byte in an isolated directory.
These checks authenticate the fixed instance of the proof.  The
unbounded number-theoretic statements and inherited logical calibration
remain the stated mathematical dependencies, and no new end-to-end
Lean or object-ZFC proof object is claimed.
